\subsection{Resultados de la encuesta a docentes}
\label{sec:resultados-docentes}

Respondieron la encuesta dos docentes (D01 y D02). D01 reporta impartir de primero a sexto grado, tener más de 20 años de experiencia y usar tecnología en clase con frecuencia; D02 imparte primer grado, tiene de 11 a 20 años de experiencia y usa tecnología a veces. La Tabla~\ref{tab:resultados-docentes} presenta sus valoraciones. Con dos respuestas, la media solo resume ambos valores y no debe interpretarse como una estimación de la opinión docente en general.

\begin{table}[H]
\centering
\caption{Valoración docente de las afirmaciones (1 = totalmente en desacuerdo, 5 = totalmente de acuerdo).}
\label{tab:resultados-docentes}
\renewcommand{\arraystretch}{1.3}
\begin{tabular}{|p{9.6cm}|c|c|c|}
\hline
\textbf{Afirmación} & \textbf{D01} & \textbf{D02} & \textbf{Media} \\ \hline
1. La explicación del proyecto fue clara. & 5 & 4 & 4.5 \\ \hline
2. El problema que atiende Graf-IA es relevante para mis alumnos. & 5 & 4 & 4.5 \\ \hline
3. Los ejercicios corresponden al nivel de primer grado. & 5 & 4 & 4.5 \\ \hline
4. Los ejercicios son coherentes con los métodos de enseñanza que utilizo. & 5 & 4 & 4.5 \\ \hline
5. Un niño de 6 a 7 años podría usar la vista del alumno con poca ayuda. & 5 & 3 & 4.0 \\ \hline
6. La letra, el tamaño de los textos y los íconos son adecuados. & 3 & 3 & 3.0 \\ \hline
7. El lienzo se parece a escribir en el cuaderno. & 5 & 4 & 4.5 \\ \hline
8. El botón «Escuchar ejercicio» es útil para alumnos que aún no leen con fluidez. & 5 & 4 & 4.5 \\ \hline
9. La calificación sugerida (ortografía y legibilidad) es fácil de entender. & 5 & 5 & 5.0 \\ \hline
10. Es adecuado que el maestro confirme o corrija la calificación antes de que cuente. & 5 & 5 & 5.0 \\ \hline
11. La vista de progreso ayudaría a identificar a los alumnos que necesitan apoyo. & 5 & 5 & 5.0 \\ \hline
12. Estaría dispuesto(a) a usar Graf-IA con mi grupo. & 5 & 4 & 4.5 \\ \hline
\end{tabular}
\end{table}

Ambos docentes otorgaron la valoración máxima a las afirmaciones sobre la claridad de la calificación sugerida, la confirmación docente y la vista de progreso (afirmaciones 9 a 11). La afirmación con menor valoración fue la adecuación de la letra, el tamaño de los textos y los íconos, a la que ambos asignaron un 3. D02 también asignó un 3 a la posibilidad de que un niño de 6 a 7 años use la vista del alumno con poca ayuda.

Sobre la frecuencia de uso prevista, D01 indicó que usaría la aplicación una vez por semana y D02 que la usaría ocasionalmente. Ambos señalaron la falta de internet y la falta de dispositivos como dificultades para usarla en su escuela; ninguno marcó la falta de tiempo en clase ni la capacitación docente. En la pregunta abierta, D01 respondió que no cambiaría nada, y D02 indicó que lo primero sería «conocer el programa, saberlo utilizar, para poder ayudar, enseñar y guiar a los alumnos».
